% Appendix B and C text drafted by role A from the rule_v1 audit (FINAL_TASKS A, P1).
% Every number is a result key (docs/RESULT_KEYS.md): sum/A-AUDIT/<field> reads results/A-AUDIT/summary.json
% (scripts/audit_rule_v1.py -> tables/data_stats.json -> scripts/render_docs.py); sum/A-SETS/<field> reads
% results/A-SETS/summary.json (scripts/export_set_summaries.py). Wording follows docs/PAPER_VS_CODE.md.
% The authors rewrite this in their own words (H6); the facts and keys must survive the rewrite.

% ---------------------------------------------------------------- Appendix B
\paragraph{Manifest and provenance.}
The released manifest gives, for each of the \res{sum/A-AUDIT/rules.library@int} rules:
\begin{itemize}
\item its source or specification;
\item the text shown in prompts;
\item the program and the applicability semantics of each criterion;
\item its tests and its version.
\end{itemize}
The \res{sum/A-AUDIT/rules.sampled_by_kind.constraint@int} grammar-sampled rules are synthetic.
For each of the \res{sum/A-AUDIT/rules.hand_written@int} hand-written rules, we searched for the source
document. We recorded the supporting text, every simplification and every contradiction with the source.
\begin{itemize}
\item The source was found and opened for \res{sum/A-AUDIT/rules.sources_verified_online@int} rules.
\item \res{sum/A-AUDIT/rules.hand_written_provenance.simplified_source@int} rules are simplified source
rules and \res{sum/A-AUDIT/rules.hand_written_provenance.synthetic@int} are synthetic; none restates its
source verbatim.
\item \res{sum/A-AUDIT/rules.sources_with_contradictions@int} rules contradict their source in at least one
respect, for example a threshold or an alternative the source does not support.
\end{itemize}
Model agents carried out this search, and a second agent tried to refute each record. An author has not
yet checked the formalisations against their sources. Because every prompt states the rule, correctness is
relative to the stated text for all rules, and no rule is presented as clinical guidance.

\paragraph{Constraint rules and claims.}
Each constraint rule states a default option, one or more conditions and an alternative. The hand-written
rules were written directly as programs together with their rule text. Tests execute every threshold one
step below, at and above it, and check that the operator wording of the text matches the program. The
remaining constraint rules are sampled from a typed grammar. The sampler never pairs a condition with an
alternative that the condition rules out, and never pairs it with a setting that the condition
contradicts. The two options of a rule are plain orders whose lengths differ by at most
\res{sum/A-AUDIT/library_semantics.constraint_claim_token_diff_max@int} tokens.

\paragraph{Rendering audit.}
Before the freeze, \res{sum/A-AUDIT/reviewers.sessions@int} LLM agent sessions read
\res{sum/A-AUDIT/reviewers.groups@int} rendered groups in two rounds, and no person took part. Each session
checked four things:
\begin{itemize}
\item that every stated answer follows from a literal reading of the rule and the case;
\item consistency and plausibility;
\item rendering;
\item give-aways.
\end{itemize}
No stored label disagreed with the program. In the first round, the sessions marked some answers as
disputable under a literal reading. These were mostly missing-input lines worded ``not recorded'', which
reads as absent, and a few condition labels and time phrases. Their wording was changed. The second round
found none. Presentation edits that changed a fact, alternatives ruled out by their condition, and
near-misses that differed from the flip in more than one attribute were removed. These sessions read
pre-freeze versions of the generator. The authors' reading of \res{sum/A-AUDIT/sample_sheets.groups@int}
frozen groups is the human check of text fidelity.

\paragraph{Splits and overlap.}
In each of \res{sum/A-AUDIT/rules.folds@int} folds, 5 of the \res{sum/A-AUDIT/rules.signature_classes@int}
signature classes are held out for L2. Between the training rules of fold 1 and L2, the following share
anything with training:
\begin{itemize}
\item canonical programs: \res{sum/A-AUDIT/overlap.L2.program_pct@1}\%;
\item parameter-normalised structures: \res{sum/A-AUDIT/overlap.L2.structure_pct@1}\%;
\item rule texts: \res{sum/A-AUDIT/overlap.L2.rule_text_pct@1}\%;
\item rendering templates: \res{sum/A-AUDIT/overlap.L2.templates_shared@int};
\item base states: \res{sum/A-AUDIT/overlap.L2.base_states_shared@int};
\item source documents, among L2 rules with a source:
\res{sum/A-AUDIT/overlap.L2.source_shared_with_train_pct@1}\%.
\end{itemize}
Criterion subexpressions are mostly shared (\res{sum/A-AUDIT/overlap.L2.subexpression_pct@1}\%), so L2
tests new compositions of familiar conditions. Cue words for relatives, negation and time are split in
the same way as templates. The line that a flip or near-miss adds never contains a training cue word
(\res{sum/A-AUDIT/overlap.L2.near_lines_with_train_cue_pct@1}\% of near-miss lines). Among rules shared
with training (L0, development, L3-alt), base states can recur:
\res{sum/A-AUDIT/overlap.L0.base_states_shared@int} of \res{sum/A-AUDIT/overlap.L0.base_states@int} in L0.
A rule has few distinct states.

\paragraph{L3-alt.}
Each of the \res{sum/A-AUDIT/l3alt.groups@int} groups moves exactly one threshold of a measured criterion.
The flip lies between the original and the altered threshold
(\res{sum/A-AUDIT/l3alt.flips_between_thresholds_pct@1}\%), so the original rule gives the other answer
on it. Its near-miss kinds are therefore the numeric ones only: numeric
(\res{sum/A-AUDIT/l3alt.near_miss_kinds.numeric@int}), at the threshold
(\res{sum/A-AUDIT/l3alt.near_miss_kinds.boundary@int}) and an older value
(\res{sum/A-AUDIT/l3alt.near_miss_kinds.time@int}).

% ---------------------------------------------------------------- Appendix C
\paragraph{Invariants.}
Every edit is checked on the target claim by executing the program. A flip must change the claim, and a
near-miss or a presentation edit must not; this holds for rules that combine criteria too. Flips are
decisive by construction, because the other criteria are held met or unmet so that the target alone
decides the conclusion. No proposal is ever rejected because a second criterion changes. The only rejected
proposals are presentation edits that would change nothing but the header
(\res{sum/A-AUDIT/rejects.logged.L2.rejected@int} of \res{sum/A-AUDIT/rejects.logged.L2.proposed@int} in
L2). A missing twin is kept only if two completions of the decisive input, the base and the flip mention,
reach both outcomes. This is sufficient for ``undetermined''.
\begin{itemize}
\item The equalities $h_k(x)=0$, $h_k(\xf)=h_k(\xn)=1$ ($h_k\equiv 1$ for measured inputs) are checked on
the rendered text, and every ledger quote must occur in its case.
\item The program runs on the generator's state; nothing is parsed from the text.
\item A re-execution of all labels, invariants and missing twins on every test and development set
(\res{sum/A-AUDIT/target_claim_check.groups_total@int} groups) found
\res{sum/A-AUDIT/target_claim_check.violations@int} violations.
\end{itemize}

\paragraph{Semantics.}
\emph{Several mentions.} A finding criterion holds if any mention it counts has status present. A
measurement criterion counts exactly one value, the patient's current value. Earlier values are reported
as past and are not counted; every case states exactly one current value of each measured input
(\res{sum/A-AUDIT/semantics.all_test_and_dev.numeric_inputs_one_counted_value@int} of
\res{sum/A-AUDIT/semantics.all_test_and_dev.numeric_inputs@int}). \emph{Conflicting values.} A case holds
at most one mention of each finding, so a counted ``absent'' and a counted ``present'' never meet.
\emph{Numbers.} Each measured input has one unit across all templates, and the rule text states it
(\res{sum/A-AUDIT/library_semantics.numeric_criteria_unit_in_rule_text@int} of
\res{sum/A-AUDIT/library_semantics.numeric_criteria_with_a_unit@int} criteria). No edit changes a unit.
Thresholds are strict or inclusive as the text says. A value exactly at a strict threshold is its own
near-miss kind; at an inclusive threshold the same value is a flip.
\emph{Time.} A mention is current or past, and
\res{sum/A-AUDIT/semantics.all_test_and_dev.past_mentions_with_year@int} of
\res{sum/A-AUDIT/semantics.all_test_and_dev.past_mentions@int} past mentions carry a year. Cases state no
reference date and rule\_v1 has no time windows. The rule-side items and the registered criteria state a
visit date; there, ``within $W$ months'' counts the day exactly $W$ months earlier, and the rule text says
so. \emph{Omission.} Measurements that are not given are unknown, and history items that are not mentioned
are absent, unless the case says they are unknown. This is a stipulation of the generator. No finding base
names the decisive concept.
\begin{itemize}
\item \res{sum/A-AUDIT/semantics.all_test_and_dev.finding_bases_generic_absence_line@int} of
\res{sum/A-AUDIT/semantics.all_test_and_dev.finding_bases@int} finding bases carry a general line implying
its absence.
\item \res{sum/A-AUDIT/semantics.all_test_and_dev.finding_bases_no_line@int} say nothing.
\end{itemize}

\paragraph{Near-miss kinds.}
In L2 each kind accounts for 20\% of the triplets: numeric (\res{sum/A-AUDIT/near_miss_kinds.L2.numeric.n@int}),
at a strict threshold (\res{sum/A-AUDIT/near_miss_kinds.L2.boundary.n@int}), another subject
(\res{sum/A-AUDIT/near_miss_kinds.L2.subject.n@int}), explicit absence
(\res{sum/A-AUDIT/near_miss_kinds.L2.negation.n@int}) and a time outside the predicate
(\res{sum/A-AUDIT/near_miss_kinds.L2.time.n@int}). For criteria whose predicate counts past or family
findings, such mentions are flips.

\paragraph{Harder cases, by requirement.}
L2 triplets per requirement (a triplet can meet several):
\begin{itemize}
\item rule-dependent time scope: \res{sum/A-AUDIT/requirements.L2.rule_dependent_time_scope@int};
\item an older value beside the current one: \res{sum/A-AUDIT/requirements.L2.temporal_selection@int};
\item composition with another condition met:
\res{sum/A-AUDIT/requirements.L2.composition_with_other_condition_met@int};
\item numerical semantics at a threshold: \res{sum/A-AUDIT/requirements.L2.numerical_semantics@int};
\item long note: \res{sum/A-AUDIT/requirements.L2.long_note@int};
\item superseded value: \res{sum/A-AUDIT/requirements.L2.superseded_value@int};
\item none of these: \res{sum/A-AUDIT/requirements.L2.none_of_these@int}.
\end{itemize}
Competing mentions of the form ``a relative has the finding; the patient explicitly does not'' are not
generated.

\paragraph{Rule-side items.}
Each of the \res{sum/A-SETS/xr_v1.n_items@int} items keeps a case fixed and pairs two rules that differ
only in the applicability clause:
\begin{itemize}
\item time window (\res{sum/A-SETS/xr_v1.by_dimension.window@int});
\item currency (\res{sum/A-SETS/xr_v1.by_dimension.currency@int});
\item subject (\res{sum/A-SETS/xr_v1.by_dimension.subject@int});
\item threshold inclusivity, upper and lower (\res{sum/A-SETS/xr_v1.by_dimension.inclusivity@int}).
\end{itemize}
There are three phrasings per clause. The three cases are the contested form, an uncontested positive and
a negative, and an item is solved only if all six judgments are right. A scorer that ignores the rule text,
one that never counts the contested form and one that always counts it all reach crossed accuracy
\res{sum/A-SETS/xr_v1.validation.crossed_accuracy.ignores_rule_text_counting_rule.XA@int}; the program
reaches \res{sum/A-SETS/xr_v1.validation.crossed_accuracy.program.XA@int}. All rules are synthetic
interventions on grammar scenarios.
